%% ==========================================================================
%% DRAFT -- August 6, 2026. Author byline: Daniel Kirtchakov (05oz,
%% halfounce.io; no institutional affiliation).
%%
%% Object re-verified 2026-08-06 by two independently written standard-library
%% verifiers (verify_cfr525.py, shipped with this note; verify_cfr.py, in the
%% program's construct/florentine directory). Both confirm all 5*25*24 = 3000
%% ordered distance events distinct; SHA-256 of CFR_5_25.json matches the
%% recorded value 9e2d9b33...779ffdcc.
%%
%% Claim strength imposed by the program's gate and honored throughout:
%%   * The rectangle establishes F_c(25) >= 5, ONE MORE than the value 4
%%     recorded for n = 25 in the Handbook of Combinatorial Designs, 2nd ed.
%%     (2006), Table 62.27 (p. 677). The improvement is stated against that
%%     RECORDED value, not against the full state of knowledge.
%%   * H.-Y. Song, Comput. Math. Appl. 39 (2000), no. 11, 31-35 could not be
%%     read in full (paywall/CAPTCHA on every open route); its Table 1 was not
%%     consulted, so NO priority over Song 2000 is claimed. See Section 5.
%%   * Novelty sweep clean 2026-08-06: no surveyed source states F_c(25) >= 5;
%%     CPro1 lists (5,25) as an open instance.
%% The primary source (Handbook chapter VI.62) was read directly; the decisive
%% passages (Def. 62.26; Table 62.27, n=25 entry) are quoted at point of use.
%% ==========================================================================
\documentclass[11pt]{amsart}

\usepackage[margin=1.1in]{geometry}
\usepackage{amsmath,amssymb}
\usepackage{booktabs}
\usepackage{array}
\usepackage[colorlinks=true,allcolors=blue]{hyperref}

\emergencystretch=3em
\tolerance=1200
\hbadness=4000

\newcommand{\Fc}{F_{c}}
\newcommand{\Zn}{\mathbb{Z}}

\theoremstyle{plain}
\newtheorem{theorem}{Theorem}[section]
\newtheorem{lemma}[theorem]{Lemma}
\newtheorem{proposition}[theorem]{Proposition}
\newtheorem{corollary}[theorem]{Corollary}
\theoremstyle{definition}
\newtheorem{definition}[theorem]{Definition}
\theoremstyle{remark}
\newtheorem{remark}[theorem]{Remark}

\begin{document}

\title[A $5\times 25$ circular Florentine rectangle: $\Fc(25)\ge 5$]{A
$5\times 25$ circular Florentine rectangle: $\Fc(25)\ge 5$, exceeding the
recorded lower bound}

\author{Daniel Kirtchakov}
\address{Independent researcher (\texttt{05oz}); no institutional
affiliation.}
\email{daniel@halfounce.io}
\urladdr{https://halfounce.io}

\thanks{\emph{Computation and authorship.} The search that produced the
object below, and this note, were carried out by \textbf{Claude Fable~5}
(Anthropic), directed by the author, on a single Apple~M4 laptop. This is a
methods statement, and it is the point of the note: the object is checked by
an exhaustive, exact standard-library verifier, so the provenance of the
\emph{search} is irrelevant to the validity of the \emph{result} (see
\S\ref{sec:tcb}). For the record, the object was located by a C program
(\texttt{cfr\_search.c}) as a maximum clique, among the $K$-equivariant
orthomorphisms of $\Zn_{25}$ for the prescribed multiplier group
$K=\langle 7\rangle=\{1,7,18,24\}\le\Zn_{25}^{\ast}$; none of that enters the
verification.}

\thanks{\emph{Prior-art record.} The primary source, chapter VI.62 (``Tuscan
Squares,'' by W.~Chu, S.~W.~Golomb and H.-Y.~Song) of the \emph{Handbook of
Combinatorial Designs}, 2nd ed.\ \cite{HCD}, was read directly on
August~6, 2026; the two passages the claim rests on --- Definition 62.26 and
the $n=25$ entry of Table 62.27 (p.~677) --- are quoted verbatim in
\S\ref{sec:claim}. The original H.-Y.~Song paper \cite{Song00} could not be
read in full (its publisher page was inaccessible on every open route tried on
that date); its Table 1 was not consulted, and no priority over \cite{Song00}
is claimed (\S\ref{sec:novelty}). A full-text novelty sweep of the same date
found no surveyed source asserting $\Fc(25)\ge 5$; the constructive benchmark
CPro1 \cite{CPro1} lists $(5,25)$ as an open instance. \emph{Draft of
August~6, 2026.}}

\subjclass[2020]{Primary 05B30; Secondary 05B15, 68V15}
\keywords{Circular Florentine rectangle, Tuscan square, orthomorphism,
exhaustive verification, certified computation}

\date{August 6, 2026}

\begin{abstract}
An $r\times n$ circular Florentine rectangle is an array of $r$ rows, each a
permutation of $\{0,1,\dots,n-1\}$ read circularly, such that for every
ordered pair of distinct symbols $(a,b)$ and every distance
$m\in\{1,\dots,n-1\}$ at most one row places $b$ exactly $m$ positions
circularly to the right of $a$; $\Fc(n)$ denotes the largest $r$ for which one
exists. We exhibit an explicit $5\times 25$ circular Florentine rectangle and
verify it exhaustively and exactly: its $5\cdot 25\cdot 24 = 3000$ ordered
distance events are pairwise distinct, checked two independent ways by a
standard-library program that shares no code with the search, and again by a
second, independently written verifier. The object establishes
$\Fc(25)\ge 5$ --- one more than the value $4$ recorded for $n=25$ in the
\emph{Handbook of Combinatorial Designs}, 2nd ed.\ (2006), Table 62.27
(p.~677), where the entry reads ``$4\ \cdots\ 24$'' and the lower bound $4$ is
exactly the $p-1$ of the multiplier construction (Construction 62.22(1),
$p=5$ the least prime factor of $25$). The trusted base contains nothing
beyond the verifier and its interpreter: there is no solver to trust, no
enumeration whose completeness is assumed, and no unproven step. We do not
claim priority over H.-Y.~Song's 2000 paper, which we could not read; the
claim is stated against the recorded Handbook value, and the array itself is
new to the surveyed literature.
\end{abstract}

\maketitle

\section{The object and the property}\label{sec:object}

Fix $n\ge 2$ and the symbol set $S=\{0,1,\dots,n-1\}$. Read the positions of a
length-$n$ row circularly, i.e.\ modulo $n$.

\begin{definition}[circular Florentine rectangle; \protect{\cite[Def.\ 62.1]{HCD}}]
\label{def:cfr}
An $r\times n$ \emph{circular Florentine rectangle} is an $r\times n$ array in
which each row is a permutation of $S$ and which has the following property:
for every ordered pair of distinct symbols $(a,b)$ and every distance
$m\in\{1,\dots,n-1\}$, there is \emph{at most one} row in which symbol $b$
occupies the position $m$ steps circularly to the right of symbol $a$.
\end{definition}

Equivalently, form for every row, every start position $i$, and every distance
$m\in\{1,\dots,n-1\}$ the ordered \emph{distance event}
\[
(a,b,m),\qquad a=R[i],\ \ b=R[(i+m)\bmod n],
\]
where $R$ is the row. Definition~\ref{def:cfr} holds if and only if all such
events are pairwise distinct. A single row of length $n$ generates
$n(n-1)$ events, all distinct within the row because the row is a permutation;
so an $r\times n$ array has the property exactly when its $r\,n(n-1)$ events
are pairwise distinct across the whole array.

\begin{figure}[ht]
\centering
\scriptsize
\setlength{\tabcolsep}{2.6pt}
\renewcommand{\arraystretch}{1.15}
\begin{tabular}{r|*{25}{r}}
\toprule
 & \multicolumn{25}{c}{\scriptsize position $0,1,\dots,24$ (read circularly)}\\
\midrule
row $0$ & 0 & 1 & 2 & 3 & 4 & 5 & 6 & 7 & 8 & 9 & 10 & 11 & 12 & 13 & 14 & 15 & 16 & 17 & 18 & 19 & 20 & 21 & 22 & 23 & 24\\
row $1$ & 0 & 7 & 14 & 21 & 3 & 10 & 17 & 24 & 6 & 13 & 20 & 2 & 9 & 16 & 23 & 5 & 12 & 19 & 1 & 8 & 15 & 22 & 4 & 11 & 18\\
row $2$ & 0 & 24 & 3 & 12 & 16 & 20 & 14 & 18 & 2 & 6 & 15 & 4 & 8 & 17 & 21 & 10 & 19 & 23 & 7 & 11 & 5 & 9 & 13 & 22 & 1\\
row $3$ & 0 & 5 & 22 & 11 & 23 & 19 & 18 & 10 & 24 & 16 & 8 & 21 & 13 & 12 & 4 & 17 & 9 & 1 & 15 & 7 & 6 & 2 & 14 & 3 & 20\\
row $4$ & 0 & 15 & 9 & 24 & 7 & 22 & 19 & 5 & 17 & 14 & 4 & 12 & 2 & 23 & 13 & 21 & 11 & 8 & 20 & 6 & 3 & 18 & 1 & 16 & 10\\
\bottomrule
\end{tabular}
\caption{The circular Florentine rectangle $\mathrm{CFR}(5,25)$. Each of the
five rows is a permutation of the symbol set $\{0,1,\dots,24\}$; the columns
are the position indices $0,1,\dots,24$, read circularly. The array satisfies
Definition~\ref{def:cfr}: for every ordered pair of distinct symbols $(a,b)$
and every distance $m\in\{1,\dots,24\}$, at most one row places $b$ exactly
$m$ positions circularly to the right of $a$; equivalently, the
$5\cdot 25\cdot 24 = 3000$ ordered distance events $(a,b,m)$ are pairwise
distinct. Rows $0$ and $1$ are the linear permutations $t\mapsto t$ and
$t\mapsto 7t\pmod{25}$; rows $2$, $3$, $4$ are non-linear. The recorded
digest of the object is
$\mathtt{sha256}=\mathtt{9e2d9b33\ldots779ffdcc}$.}
\label{fig:cfr}
\end{figure}

\section{The claim}\label{sec:claim}

\subsection*{What $\Fc$ records}
The primary source defines the quantity our object bounds. Quoting
\cite[Def.\ 62.26, p.~677]{HCD} verbatim:
\begin{quote}
``Let $\Fc(n)$ denote the maximum integer $\Fc$ such that an
$\Fc\times n$ circular florentine rectangle exists.''
\end{quote}
The basic bounds recorded there are $p-1\le \Fc(n)\le n-1$, where $p$ is the
least prime factor of $n$: the lower bound is the multiplier construction
\cite[Construction 62.22(1), p.~676]{HCD}, which for the least prime factor
$p$ of $n$ takes the rows $t\mapsto ct\pmod n$ over a suitable set of
multipliers $c$ and yields a $(p-1)\times n$ circular Florentine rectangle;
the upper bound $n-1$ is immediate. For $n=25$ the least prime factor is
$p=5$, so the construction gives $4$ rows, and \cite[Table 62.27,
p.~677]{HCD} --- headed ``Possible values of $\Fc(n)$ and $F(n)$'' --- records
the $n=25$ entry, verbatim, as
\[
25\qquad 4\ \cdots\ 24 ,
\]
i.e.\ a best recorded lower bound of $4$ (the value $p-1$) and the trivial
upper bound $24=n-1$; the value of $\Fc(25)$ within that range is left
undetermined.

\subsection*{The result}

\begin{theorem}\label{thm:main}
The array of Figure~\ref{fig:cfr} is a circular Florentine rectangle. Hence
\[
\boxed{\ \Fc(25)\ \ge\ 5\ }
\]
which is one more than the value $4$ recorded for $n=25$ in the
\emph{Handbook of Combinatorial Designs}, 2nd ed.\ (2006), Table 62.27
(p.~677).
\end{theorem}

\begin{proof}
By Definition~\ref{def:cfr} it suffices that the $5\cdot 25\cdot 24 = 3000$
ordered distance events of the five rows in Figure~\ref{fig:cfr} are pairwise
distinct; this is confirmed exhaustively in \S\ref{sec:verif}. The array is
then a $5\times 25$ circular Florentine rectangle, so by
\cite[Def.\ 62.26]{HCD} the maximum $r$ for which an $r\times 25$ such
rectangle exists is at least $5$.
\end{proof}

The two linear rows $t\mapsto t$ and $t\mapsto 7t$ are multiplier rows of the
form $t\mapsto ct$ that underlies the multiplier construction
\cite[62.22(1)]{HCD}; the remaining three rows are non-linear (this is confirmed by the verifier, which finds no
multiplier $c$ with $R[t]\equiv ct$ for rows $2$, $3$, $4$). The recorded
lower bound $4=p-1$ is exactly what the multiplier construction alone
delivers; Figure~\ref{fig:cfr} exceeds it by adjoining a fifth, non-linear
row that remains event-disjoint from the other four.

\section{Verification}\label{sec:verif}

The verification is \emph{exhaustive} and \emph{exact}: it enumerates all
$3000$ ordered distance events directly, in integer arithmetic, and confirms
they are distinct. There is no sampling, no floating point, and no randomized
or heuristic step.

\subsection*{The shipped verifier}
\texttt{verify\_cfr525.py} ($173$ lines, Python standard library only ---
\texttt{json}, \texttt{hashlib}, \texttt{sys}, \texttt{os} --- with no
POSIX-only call, so it runs unchanged on any platform with CPython~3) performs,
against the object file \texttt{CFR\_5\_25.json}:
\begin{enumerate}
\item \emph{SHA-256.} It recomputes the digest of the object file and requires
it to equal the recorded value
\texttt{9e2d9b33\ldots779ffdcc}; and it cross-checks the rows read from disk
against a second, independent in-source transcription of the same data.
\item \emph{Shape and permutation.} It checks $r=5$, $n=25$, and that each row
is a permutation of $\{0,\dots,24\}$.
\item \emph{Property, method A (literal circular window).} For every row,
every start position $i$, and every distance $m\in\{1,\dots,24\}$ it forms the
event $(a,b,m)$ with $a=R[i]$, $b=R[(i+m)\bmod 25]$, and counts the distinct
events. The property holds iff the count equals $r\,n(n-1)=3000$.
\item \emph{Property, method B (per ordered pair via positions).} For each
ordered pair $(a,b)$ of distinct symbols it computes, in each row, the unique
distance $(\mathrm{pos}[b]-\mathrm{pos}[a])\bmod 25$, and requires these to be
pairwise distinct across the rows. It then requires methods A and B to agree.
\end{enumerate}
Run on the shipped object today, the verifier reports, in its final
lines (the long verdict line soft-wrapped here to fit the text block),
\begin{verbatim}
[sha256] match    = True
[shape ] r=5, n=25, 5 rows of length n: True
[cross ] disk rows == independent transcription: True
[perm  ] every row is a permutation of {0..24}: True
[meth A] window events 3000/3000 distinct (expected 3000): True
[meth B] per-pair distances 3000/3000 distinct (expected 3000): True
[cross ] methods A and B agree: True

VERIFIED: the object is a valid CFR(5,25); all 3000 ordered distance
events are distinct, so F_c(25) >= 5.
\end{verbatim}
As a negative control, swapping any two symbols within a row (which breaks the
digest and the property) is rejected: the event count falls below $3000$ and
the run exits non-zero.

\subsection*{A second, independent verifier}
A second verifier, \texttt{verify\_cfr.py} (in the program's
\texttt{construct/florentine} directory), was written separately from the
first and from the search. It re-checks the same object through unrelated code
--- the SHA-256, the permutation property, both event-counting methods, and,
additionally, the $K=\langle 7\rangle$ multiplier equivariance and the per-row
linearity classification --- and returns the same verdict
(\texttt{OVERALL: OBJECT IS A VALID CFR(5,25): True}). The two verifiers share
no code and agree.

\section{The trusted base}\label{sec:tcb}

What a skeptic must believe to accept Theorem~\ref{thm:main} is unusually
small, and it is worth stating exactly, because it is smaller than in the
program's earlier certified notes (each of which declared at least one trust
assumption, such as the completeness of an external enumerator).

\emph{Machine-checked, from raw data:} that the object file has the recorded
SHA-256; that each row is a permutation of $\{0,\dots,24\}$; and that all
$3000$ ordered distance events are pairwise distinct, established two
independent ways and reproduced by a second, independently written verifier.

\emph{Assumed:} only that CPython and the operating system execute the
$173$-line \texttt{verify\_cfr525.py} as written, and that its transcription of
Definition~\ref{def:cfr} (\cite[Def.\ 62.1/62.26]{HCD}) into code is faithful
--- which a reader confirms by reading the verifier against the definition.
There is nothing else. There is no external solver, no enumeration whose
completeness must be trusted, no unproven lemma, and no floating-point or
randomized computation. The check is a finite, exact, exhaustive enumeration
of $3000$ events, and it either finds all of them distinct or it does not.

In particular, the search apparatus --- \texttt{cfr\_search.c}, the multiplier
group $K$, the clique computation --- is \emph{not} in the trusted base. The
object stands or falls on its own, independently of how it was found.

\section{Novelty, and what is not claimed}\label{sec:novelty}

\emph{The improvement is over the recorded value.} Theorem~\ref{thm:main}
asserts $\Fc(25)\ge 5$ and states the improvement against the value $4$
recorded for $n=25$ in Table 62.27 (p.~677) of \cite{HCD} --- a citable
secondary source whose $n=25$ entry reads ``$4\ \cdots\ 24$'' and which cites,
for that table, references predating the specific determination of $\Fc(25)$.
A full-text novelty sweep on August~6, 2026 found no surveyed source asserting
$\Fc(25)\ge 5$: the recent literature that uses circular Florentine rectangles
quotes only the basic bounds $p-1\le\Fc(N)\le N-1$, and the constructive
benchmark CPro1 \cite{CPro1} lists $(5,25)$ among its open instances, with a
property definition matching Definition~\ref{def:cfr}. The array of
Figure~\ref{fig:cfr} is new to the surveyed literature.

\emph{No priority over Song 2000 is claimed.} The one paper that might already
contain a determination of $\Fc(25)$ is H.-Y.~Song, \emph{The existence of
circular Florentine arrays} \cite{Song00}. Its publisher page was inaccessible
on every open route tried on August~6, 2026, so its Table 1 was not consulted.
We therefore make \emph{no} claim of priority over \cite{Song00}. The claim of
Theorem~\ref{thm:main} is deliberately anchored on the recorded Handbook value
of $4$, not on the assertion that $4$ was the best value obtainable anywhere:
we state that the object exceeds the recorded lower bound, and no more.

\section{Artifacts}\label{sec:artifacts}

Two files are shipped with this note; both are pinned by SHA-256.
\begin{center}
\small
\begin{tabular}{lll}
\toprule
file & bytes & SHA-256\\
\midrule
\texttt{CFR\_5\_25.json} & $1236$ &
\texttt{9e2d9b33\ldots779ffdcc}\\
\texttt{verify\_cfr525.py} & --- &
\texttt{9d7363d0\ldots37dd0bb}\\
\bottomrule
\end{tabular}
\end{center}
\noindent The full digests are
\begin{quote}\ttfamily\footnotesize
CFR\_5\_25.json\ \ \ \ 9e2d9b339a79d13783d2e24efde49e5ed7904418f2179bf7d3f0143d779ffdcc\\
verify\_cfr525.py\ \ 9d7363d08ce2cd90b194ce1f8c318545501006c14b3ea1cb7871586ec37dd0bb
\end{quote}
To reproduce the verification on any machine with CPython~3, from the
directory holding the two files:
\begin{verbatim}
python3 verify_cfr525.py          # reads CFR_5_25.json beside the script
python3 verify_cfr525.py PATH     # or an explicit path to the object
\end{verbatim}
No compiler, solver, or third-party package is required; the script imports
only the Python standard library and terminates in well under a second.

\section*{Acknowledgments}

The definitions, the quantity $\Fc(n)$, the multiplier construction that fixes
the recorded lower bound $4$, and the value tables are those of the
\emph{Handbook of Combinatorial Designs} chapter on Tuscan squares by Wai~Chu,
Solomon~W.~Golomb and Hong-Yeop~Song \cite{HCD}. The $(5,25)$ instance is
posed as open by the CPro1 benchmark \cite{CPro1}. The computation and drafting
were AI-assisted as stated in the first footnote; the object is checked by an
exhaustive standard-library verifier, so that assistance bears on how the
object was found, not on whether it is valid.

\begin{thebibliography}{9}

\bibitem[HCD]{HCD}
C.~J.~Colbourn and J.~H.~Dinitz (eds.),
\emph{Handbook of Combinatorial Designs}, 2nd ed.,
CRC Press, Boca Raton, 2006. Chapter VI.62, ``Tuscan Squares,'' by
W.~Chu, S.~W.~Golomb and H.-Y.~Song, pp.~673--678. Read directly on
August~6, 2026; the claim rests on Definition 62.26 and the $n=25$ entry of
Table 62.27 (p.~677), ``$4\ \cdots\ 24$,'' with the lower bound $4=p-1$ from
Construction 62.22(1) ($p=5$).

\bibitem[Son00]{Song00}
H.-Y.~Song,
\emph{The existence of circular Florentine arrays},
Comput.\ Math.\ Appl.\ \textbf{39} (2000), no.~11, 31--35.
DOI 10.1016/S0898-1221(00)00104-8. Not read in full (publisher page
inaccessible on every open route tried on August~6, 2026); its Table~1 was not
consulted, and no priority over this paper is claimed.

\bibitem[CPro1]{CPro1}
Constructive-Codes, \emph{CPro1} benchmark suite,
constructive problem \texttt{circular-florentine-rectangle}, which lists the
open instances $(6,21)$, $(5,25)$, $(5,27)$, $(4,33)$ with the property of
Definition~\ref{def:cfr}. On-disk copy consulted August~6, 2026.

\end{thebibliography}

\end{document}
